\chapter{62}

\section{Moosfluh - Bettmeralp (CH), Mittwoch, 29.
August}



Der Tag begann und endete als Hirtentag. Die Schafe hielten sich bis gegen Abend versteckt. Wie erwartet. Einzig schwarze Flecken,
zwischen den Unterschlüpfen vereinzelt als bewegliche, dunkle Punkte
sichtbar, führten Louis zu ihren Rückzugsorten. Mit jedem Kontakt,
jeder Berührung, mit jedem Blick in ihre dunklen Fellgesichter empfand Louis
die Begegnungen vertrauter. 

Gebannt sass er gegen Abend in ihrer Nähe und wartete. Einem Schauspiel
gleich schienen sie zu erwachen. Ein Schaf nach dem andern trottete gleichmütig auf die Ebene hinaus. 
Die Tiere verteilten sich über den Raum. Bald hüpften sie herum und schüttelten sich, als wollten sie ihr
Haarkleid lüften. Es fühlte sich leicht an, hier zu sein. Und das Treffen nach
zwölf mit Mazi hatte dem Tag die Mitte gegeben.

Rückblickend dachte Louis, waren es die Stunden dieses Tages, er,
die meiste Zeit alleine mit den Tieren inmitten dieser puren Natur, die ihn vollends
hatten ankommen lassen. Ob als Zusenn, Schafhirte oder alpiner Städter
war einerlei.

\sectionbreak

Er hatte sich nach dem Eindunkeln zur Hütte aufgemacht. 

Mit dem Reistopf auf den Knien sass er schliesslich auf der Bank vor
dem Holzhaus. Noch nie hatte ihn so viel Ruhe umgeben. Behutsam schloss er die Augen. 
Es war geradezu abenteuerlich, sich ins nächtliche Nichts hineinzuhorchen. So musste sich die Ewigkeit anfühlen.
Jeder Bissen schmeckte und nachdem er auch die letzten Reiskörner zusammengekratzt hatte, 
stellte er den Topf im Küchenbereich auf das kleine Tischchen. Aufräumen und sauber machen konnte bis morgen warten. Er würde die Hütte in Ordnung bringen. Wie er es Joh zugesichert hatte.

Wäre nicht kurz vor dem Einschlafen die Störung angeschlichen gekommen,
nicht aufzuhalten und eindringlich, hätte sich Louis richtig gut gefühlt.
Doch Valentina war zurück. Ihr böses Gesicht der Nacht starrte ihm
entgegen. Louis riss die Augen auf und fürchtete sich davor, sie wieder
zu schliessen, fürchtete sich vor dem nächsten Albtraum. 
Louis stand auf, griff sich Handy und Kopfhörer und stieg die Leiter 
hinunter. Vor der Hütte setzte er sich auf die schmale Holztreppe.
Irgend etwas musste er tun. Gedankenversunken scrollte er auf Spotify durch seine neue Playlist.
Nichts passte.
Louis stutzte. Mit welchem Songtipp hatte ihn Valentina vor Wochen aufmuntern wollen? 
Sie hatten in der Bar im Centro gesessen. Er hatte ihr von Giorgias Trennungswunsch erzählt, war
hilflos, wütend und traurig gewesen. Alles auf einmal. Sie hatte ihm stumm zugehört. Bevor sie sich verabschiedeten, blickte sie ihn eindringlich an. Die Worte waren zurück.
«Weisst du, immer, wenn ich mich mies fühle,» hatte sie gesagt, «höre ich mir \emph{«Unstoppable»} von Sia an.»
Sie hatte sanft gelächelt, ihn umarmt und war gegangen. Louis hatte ihre Worte registriert. Mehr nicht. Der Song wurde zwar immer mal wieder irgendwo gespielt. Wirklich hingehört hatte er nie. Bis jetzt. Er suchte den Titel und steckte sich die Kopfhörer ein.

\qrblock{Sia \emph{«Unstoppable»}}{media/QR-Codes/039_Unstoppable_Sia_Spotify.png}

Mittendrin stand Louis auf. Er ging auf und ab, bis er sich nach dem
zweiten Durchlauf oberhalb der Hütte wiederfand. Der Song wirkte belebend.
\emph{«Ich bin ein Porsche ohne Grenzen»}. Er blieb nachdenklich stehen.
Die Worte holten Kinderszenen hoch. Sein Versuch, sie zu ignorieren, misslang. Schon tauchte der Kleine vor ihm auf. 
Spätestens als Schüler hatte sich sein Gefühl der Einzigartigkeit immer häufiger in Selbstzweifel verkehrt. Wie ungeduldig Vater wurde, wenn Louis seine Erwartungen nicht erfüllte, wenn er doch nicht der Schnellste und Beste war. Wie hatte sich Louis vor seinem enttäuschten, oder vielmehr ärgerlichen, Gesichtsausdruck gefürchtet. Sich geschämt, versagt zu haben. Und Louis sah dieses fadenscheinige Konstrukt mit einer Klarheit vor sich, die ihn erschauern liess. Vater hatte es permanent und minutiös aufgebaut und konsequent gewartet. Dennoch hatte Louis lange und vehement versucht, ihn gross zu halten.

Es war an einem der letzten Tage gewesen, bevor sein Vater endgültig gegangen war. 
Louis hatte ein Gespräch zwischen Maman und Grandpère belauscht. 
Sie hatten ihn nicht in die Wohnung hochkommen hören, wähnten ihn wohl draussen am Spielen.
Louis’ Zimmertür hatte offen gestanden, Maman gedämpft gesprochen. Ihre Stimme hatte diesen besorgten Unterton.

«Ich bin wirklich verzweifelt, Papa. Er benutzt Louis, missbraucht ihn
und seine immer wieder komplett neuen Einfälle jagen sich. Carlos ist
ein Psycho. Mir bleibt jedes Mal die Luft weg.»

Beim Namen des Vaters war Louis zusammengezuckt. Dann hatte er
Grandpères Stimme gehört.

«Komm her, Josephine,» und nach einer kleinen Pause, «ich weiss, es ist
schwer. Du gehst grossartig damit um. Es ist richtig, dass du den
Kindern gegenüber loyal bleibst. Ich bewundere dich. Die Dinge werden
sich weisen, die Zeit arbeitet für dich.»

Dann war Mamans leises Schluchzen durch den Flur bis zu ihm
durchgedrungen. Louis hatte sich schlecht gefühlt, hatte zwar nicht
jedes Wort verstanden, doch den Ausdruck \emph{Psycho} kannte er vom Pausenhof. So
wurde der hinterhältigste Junge der Schule geschimpft. Wie konnten sie
nur so gemein von Vater sprechen. Seine Schwester hatte das Gespräch beendet, indem sie ins Wohnzimmer gestürmt war.

\sectionbreak

Louis erhob sich nachdenklich. Er staunte, wie präzise er das Gespräch der beiden gespeichert hatte. 
Unweigerlich kehrte auch das damit verbundene Gefühl der Ohnmacht zurück. Und ihm war, als stünde er gerade hilflos hinter der offenen Tür seines Kinderzimmers.

Valentina, der Gedanke kam als Rettung, sie würde ihn verstehen.
Mit welch unfassbarem Tempo sich seine Stimmungen jagten. Louis griff sich ins Haar. Er wollte
sie anrufen. Und es musste jetzt sein. Hoffentlich war sie noch wach.

«Jetzt.»

Mit festen Schritten begab er sich auf die Westseite der Hütte und
setzte sich unter den offenen Himmel. Der Ort hatte sich als
empfangssicher erwiesen. Ihre Nummer? Louis verzerrte das Gesicht. Über
die drei letzten Fünf hatte sie immer mal wieder geheimnisvoll
gelächelt. Die restlichen Zahlen fielen ihm nicht ein. Vielleicht hatte sie ihre Handynummer
auf Facebook angegeben? Bevor er \emph{Infos} anklickte, presste er die Augen
zusammen. Bitte, Valentina, bitte. Tatsächlich. Da war sie und vor ihm
lag die Zahlenreihe mit drei Mal der fünf am Ende. Er tippte die Nummer
ein.
